\subsection{A length bound for the four rational characters of order five}

\begin{proposition}[KT-ORDER-FIVE-LENGTH-ONE]
Let $p$ be prime with $p\equiv1\pmod {10}$, let
$L=\mathbf Q(\mu_p)$, and let $F$ be its unique cyclic quintic
subfield.  Write $A_p=\operatorname{Cl}(L)\{p\}$ and let $\omega$
be the Teichm\"uller character.  Then
\[
 \sum_{j=1}^4
 \operatorname{length}_{\mathbf Z_p}
 A_p^{[j(p-1)/5]}\le1.                         \tag{O5.1}
\]
Equivalently, at most one of these four components can be nonzero,
and any nonzero component is isomorphic to $\mathbf Z/p\mathbf Z$.
For every $p\ge2^{31}$ one has the explicit ordinary class-number bound
\[
                         h(F)<\frac{p^2}{6}.   \tag{O5.2}
\]
The proposition gives a size bound, not the vanishing of a specified
order-five component.
\end{proposition}

\begin{proof}
For $p<2^{31}$, the carried Hart--Harvey--Ong verification gives
$A_p^+=0$.  All four characters in (O5.1) are even, because
$(p-1)/5$ is even, so (O5.1) follows in this range.
Assume henceforth $p\ge2^{31}$.

Let $R$ be the ordinary regulator of $F$ and put $h=h(F)$.
The conductor--discriminant formula gives $d_F=p^4$.
Schoof--Washington, Corollary~(2.7), proves for every cyclic
quintic field of conductor $f$ that
\[
                        R\ge\frac1{25}\log^4(f/2).
                                                        \tag{O5.3}
\]
This is a general-field theorem in their Section~2, preceding
and independent of the Lehmer family studied in Section~3.
In particular it applies here with $f=p$.

Let $\chi$ generate the group of Dirichlet characters defining $F$.
Each of $\chi,\chi^2,\chi^3,\chi^4$ is primitive, nontrivial,
even, and of conductor $p$.  The zeta factorization and ordinary
analytic class-number formula give
\[
       \frac{16hR}{p^2}
         =\prod_{j=1}^4 L(1,\chi^j)
         =|L(1,\chi)|^2|L(1,\chi^2)|^2.         \tag{O5.4}
\]
Louboutin's even-character bound, Theorem~1(1), equation~(3), gives
$|L(1,\chi^j)|\le(\log p+\kappa_0)/2$ with
$\kappa_0=2+\gamma-\log(4\pi)<1$.
Combining this with (O5.3)--(O5.4), and writing $t=\log p>20$, yields
\[
 \begin{aligned}
 h&\le\frac{25p^2}{256}
       \left(\frac{\log p+\kappa_0}{\log(p/2)}\right)^4\\
  &<\frac{25p^2}{256}
       \left(\frac{t+1}{t-1}\right)^4
   <\frac{25\cdot21^4}{256\cdot19^4}p^2
   <\frac{p^2}{6}.
 \end{aligned}                                            \tag{O5.5}
\]
Here $t>20$ follows from $p\ge2^{31}$ and $\log2>2/3$;
we used $\log2<1$, $\kappa_0<1$, and the decreasing function
$(t+1)/(t-1)$.  The final comparison is the exact integer inequality
$6\cdot25\cdot21^4=29172150<33362176=256\cdot19^4$.
Consequently $v_p(h)\le1$.

For completeness, this bounds exactly the components in (O5.1).
Put $H=\operatorname{Gal}(L/F)$, so $p\nmid |H|$.
Extension $j$ and norm $N$ on ideal classes satisfy
\[
 Nj=|H|\quad\hbox{on }\operatorname{Cl}(F)\{p\},
 \qquad
 jN=\sum_{g\in H}g\quad\hbox{on }A_p.
                                                        \tag{O5.6}
\]
Thus $j$ identifies $\operatorname{Cl}(F)\{p\}$ with $A_p^H$:
it is injective, and for $x\in A_p^H$ one has
$x=j(|H|^{-1}N(x))$.
The characters occurring in $A_p^H$ are precisely the trivial
character and the four characters in (O5.1).
The trivial component is zero by the same norm identity for
$L/\mathbf Q$, since $\operatorname{Cl}(\mathbf Q)=0$.
The Teichm\"uller idempotents are integral over $\mathbf Z_p$;
hence their four lengths add to
$\operatorname{length}_{\mathbf Z_p}\operatorname{Cl}(F)\{p\}
=v_p(h)\le1$.  This proves (O5.1).
\end{proof}

\paragraph{Exact primary input and limitation.}
R.~Schoof and L.~C.~Washington, \emph{Quintic polynomials and real
cyclotomic fields with large class numbers},
\emph{Math. Comp.}~50 (1988), 543--556:
Section~2 begins on printed p.~544; Theorem~(2.2) and its full proof
are on pp.~545--546, Corollary~(2.6) and its proof are on p.~546,
and Corollary~(2.7) with its full proof is on p.~547.
The last corollary is precisely (O5.3).
The author-hosted primary PDF is
\url{https://reneschoof.github.io/sw.pdf}.
The proof of Corollary~(2.7) imports Pohst's five-variable product
bound $M=16$; no fresh proof audit of that earlier input is claimed.
Louboutin's exact even-character theorem and the carried
Hart--Harvey--Ong verification have the same scopes as in
KT-ORDER-SIX-COMPONENT-REDUCTION.

The quintic source explicitly states on printed p.~544 that the
quintic field of prime conductor $641491$ has class number $1566401$.
Thus an analogue of the bound $h(F)<p$ is false in degree five.
That published numerical assertion is cited as source scope; no
independent class-group recomputation is claimed here.
Neither the bound (O5.2) nor this example decides whether
$p$ can divide $h(F)$.  The contribution is the exact length-one
restriction on a possible remaining rational obstruction.
No integral $K$-group or Selmer-group length comparison is asserted.
